% Source pins (SHA-256): PifoStatement.lean=fe9475ebfa2c462d27cb0a99781926074e1281396f2c9e8d32d6dfca788d814e; pifo-noinduction-codex2.md=97eef33723fdd15eff0003d02ce1cd387a1b1dcd86a42f89302421f0fa614933; pifo-noinduction-codex.md=6aea2cf52bbeb98c855054543edc53531d6b458509b27e176f8690bec97e0a25; pifo-noinduction-fable.md=c90151cf07b362c8f49068920f68dba03fb6ca99723354d717e186ecdb068cd2.
\section{Eliminating alphabet induction and infinite alphabets}
\label{sec:infinite-alphabets}

The proof in \autoref{sec:induction} uses strong induction on the number of
packet types.  That dependency is eliminable.  In this section we give a
second organization in which the main theorem is proved simultaneously for
all finite alphabets by well-founded recursion on the two supplied finite
topologies.  The argument still uses ordinary structural induction to prove a
single-scheduler promotion lemma; what disappears is every appeal to the main
theorem merely because a type set is smaller.

The new ingredient is that a flush decomposition may be promoted all the way
to an online root interface.  A direct induction on tree height would not
establish this: the target partition can cross the actual child partition, so
its rank cannot simply be installed at the actual root.  At each internal
node we instead pass through the meet of the actual and target partitions.

\subsection{Reduced restrictions and selector laws}

For a nonempty set of types \(B\), let \(M[B]\) be the \emph{reduced
restriction} of \(M\) to \(B\).  Begin with the literal restriction, delete
inactive children, reindex the remaining child values, and recursively
suppress every internal node having only one active child.  These operations
preserve every interleaved trace.  In particular, a unary active root contains
only one child value, so each successful root pop makes exactly one anonymous
call to that child; \lemref{lem:balance} supplies the same success and failure
pattern.  Write \(\nu(M)\) for the number of nodes in this reduced topology.

It is useful to make the root construction explicit.  Let
\begin{equation}
 \mathcal C(P,r,(M_B)_{B\in P})
 \label{eq:selector-construction}
\end{equation}
denote an \(r\)-ranked root whose child selected by a type is its \(P\)-block
and whose scheduler below block \(B\) is \(M_B\).  If reduced \(M\) has an
internal root, then
\begin{equation}
 M=\mathcal C(P,r,(M_B)_{B\in P})
 \label{eq:reduced-root-presentation}
\end{equation}
up to the harmless reductions above; its actual root pair \((P,r)\)
decomposes \(F_M\).  If \(\varnothing\ne D\subseteq B\), then
\(M[D]=M_B[D]\), and this restriction has strictly smaller reduced topology
than \(M\).

We will use two consequences of the operational lemmas already proved.  First,
interleaved-equivalent corresponding children may be substituted below an
unchanged root.  The root supplies the same projected pushes and anonymous
child-pop calls, and \lemref{lem:compression} accounts for their sparse global
timestamps.  Second, if \(K\) refines \(P\), then equal selector levels
collapse:
\begin{equation}
 \mathcal C\bigl(P,u,
   (\mathcal C(K|_B,u|_B,
      (X_E)_{E\in K,\ E\subseteq B}))_{B\in P}\bigr)
 \intereq
 \mathcal C(K,u,(X_E)_{E\in K}).
 \label{eq:selector-tandem-law}
\end{equation}
Indeed, ghost-tag the references by source occurrence.  If the pending set of
the outer selector is \(\Omega\), the pending set of its inner selector below
\(B\) is \(\restrict{\Omega}{B}\), while the flat selector also has
\(\Omega\).  The outer \(u\)-minimum occurrence in \(B\) is the minimum of
\(\restrict{\Omega}{B}\); all three selectors therefore remove the same
occurrence, including when ranks tie, and invoke the same \(K\)-child.  This
is an exact occurrence coupling only because the two selector levels were
deliberately given identical \((u,\arr)\) keys.

\subsection{A finite meet that is itself a decomposition}

Recall that \((P,r)\) decomposes \(F\) when block filtering commutes with \(F\)
and the root-block word is the stable \(r\)-sorted block word, as in
\autoref{eq:block-projection} and \autoref{eq:root-pattern}.  The meet construction used in the
main proof has the following strengthened conclusion.

\begin{lemma}[Finite meet decomposition]
\label{lem:finite-meet-decomposition}
Let \(A\) be finite.  If \((P,r)\) and \((Q,s)\) decompose the same transformer
\(F:A^*\to A^*\), put
\[
 R=P\wedge Q
   =\{B\cap C:B\in P,\ C\in Q,\ B\cap C\ne\varnothing\}.
\]
There is a natural-valued total preorder \(t\) such that \((R,t)\) decomposes
\(F\), and, with equality included among the three comparison outcomes,
\begin{align}
 \cmp_t(x,y)&=\cmp_r(x,y)
   &&\text{if \(x,y\) lie in distinct \(P\)-blocks},\notag\\
 \cmp_t(x,y)&=\cmp_s(x,y)
   &&\text{if \(x,y\) lie in distinct \(Q\)-blocks}.
 \label{eq:finite-meet-interfaces}
\end{align}
The lemma uses no instance of flush-implies-interleaved equivalence.
\end{lemma}

\begin{proof}
For a cell \(D=B\cap C\), successive applications of block autonomy through
\(P\) and \(Q\) give
\begin{equation}
 \restrict{F(w)}{D}=F(\restrict{w}{D}).
 \label{eq:finite-meet-cell-autonomy}
\end{equation}
For types in distinct meet cells, prescribe the effective binary comparison
recovered by \lemref{lem:two-type}.  If \(P\) separates a pair, the
prescription is its \(r\)-comparison; if \(Q\) separates it, the prescription
is its \(s\)-comparison.  When both apply they agree because both are recovered
from the same two binary flushes.

The three-cell diamond in \lemref{lem:diamond} shows that, for any choice of
one representative from each of three distinct cells, these comparisons form
a total preorder.  Equivalently, put
oppositely directed weak edges for a prescribed tie and a marked directed
edge for a prescribed strict comparison.  The fixed-head-chord argument of
\lemref{lem:meet-preorder} excludes every directed cycle containing a marked
edge.  After collapsing the strongly connected components, a
topological ordering of the finite quotient realizes all prescriptions by
natural ranks.  This gives \(t\) and
\autoref{eq:finite-meet-interfaces}.

It remains to record the strengthening beyond
\lemref{lem:meet-preorder}.  Colored-PIFO congruence first compares \(r\) with
\(t\), using \(P\)-blocks as colors, and yields
\[
 \proj{P}(F(w))=\proj{P}(\sort_t(w)).
\]
Within a fixed \(P\)-block \(B\), apply its autonomy and compare \(s\) with
\(t\), now using \(Q\)-blocks as colors.  This yields equality of the
\(Q\)-subword inside each \(B\) between \(F(w)\) and \(\sort_t(w)\).  The
global \(P\)-word together with these per-\(P\)-block \(Q\)-subwords uniquely
reconstructs the sequence of intersections \(B\cap C\).  Hence
\begin{equation}
 \proj{R}(F(w))=\proj{R}(\sort_t(w)).
 \label{eq:finite-meet-root-pattern}
\end{equation}
Equations \eqref{eq:finite-meet-cell-autonomy} and
\eqref{eq:finite-meet-root-pattern} say exactly that \((R,t)\) decomposes
\(F\).  Unary partitions are harmless degenerate cases: their meet is the
other partition and its preorder may be used directly.
\end{proof}

\subsection{Promoting one supplied decomposition}

The next lemma is the point at which alphabet-size induction is avoided.

\begin{lemma}[Single-scheduler meet promotion]
\label{lem:meet-promotion}
Let \(M\) be a valid stateless scheduler on a nonempty finite alphabet, and
suppose that \((R,t)\) decomposes \(F_M\).  Then
\begin{equation}
 M\intereq \operatorname{Ref}(M;R,t)
 :=\mathcal C(R,t,(M[D])_{D\in R}).
 \label{eq:meet-promotion}
\end{equation}
This is a statement about one scheduler and one supplied decomposition; it
does not assume the main theorem on a smaller alphabet.
\end{lemma}

\begin{proof}
Proceed by well-founded induction on the reduced topology of \(M\), with the
induction statement quantified over every finite alphabet and every supplied
decomposition.

Suppose first that \(M\) is a leaf with type preorder \(\lambda\).  Its flush
transformer is \(\sort_\lambda\).  Applying the root-pattern equation for
\((R,t)\) to the words \(xy\) and \(yx\) shows that, across distinct
\(R\)-blocks,
\begin{equation}
 \cmp_t(x,y)=\cmp_\lambda(x,y),
 \label{eq:leaf-promotion-interface}
\end{equation}
including the rank-tie case.  Color the original \(\lambda\)-queue and the
new \(t\)-root by \(R\).  By \lemref{lem:colored} they select the same block
at every pop.  Below a selected block \(D\), maintain that the restricted leaf
\(M[D]\) contains exactly the \(D\)-projection of the original leaf queue.
The original global \(\lambda\)-minimum is also the minimum of that
projection, so the two leaves remove and emit the same occurrence.  The
residue of the \(t\)-root need not equal the residue of the original leaf;
only its block trace is coupled.

Now let \(M\) have the actual root presentation
\begin{equation}
 M=\mathcal C(P,r,(M_B)_{B\in P}).
 \label{eq:promotion-actual-root}
\end{equation}
Apply \lemref{lem:finite-meet-decomposition} to the actual decomposition
\((P,r)\) and the target decomposition \((R,t)\).  It supplies
\begin{equation}
 K=P\wedge R,\qquad u,
 \label{eq:promotion-common-meet}
\end{equation}
where \((K,u)\) decomposes \(F_M\), \(u\) agrees with \(r\) across distinct
\(P\)-blocks, and \(u\) agrees with \(t\) across distinct \(R\)-blocks.  Put
\begin{equation}
 W=\mathcal C(K,u,(M[E])_{E\in K}).
 \label{eq:promotion-W}
\end{equation}

We first compare \(M\) with \(W\).  For every \(B\in P\), restrict the two
decomposition equations for \((K,u)\) to words over \(B\).  They say precisely
that
\begin{equation}
 (K|_B,u|_B)\text{ decomposes }F_{M_B}.
 \label{eq:promotion-child-decomposition}
\end{equation}
Because \(M_B\) is an actual child, its reduced topology is strictly smaller
than that of \(M\).  The structural induction hypothesis, applied only here,
therefore gives
\begin{equation}
 M_B\intereq
 \mathcal C(K|_B,u|_B,(M[E])_{E\in K,\ E\subseteq B}).
 \label{eq:promotion-child-recursion}
\end{equation}
Substitute these schedulers below the \(P\)-root.  Change the outer root from
\(r\) to \(u\) by \lemref{lem:colored}, coloring by \(P\); only the sequence of
\(P\)-services is asserted equal.  Both selector levels now carry the same
preorder \(u\), so \autoref{eq:selector-tandem-law} collapses them to \(W\).
Thus
\begin{equation}
 M\intereq W.
 \label{eq:promotion-left-W}
\end{equation}

For the other comparison, fix \(D\in R\) and define
\begin{align}
 K_D&=\{B\cap D:B\in P,\ B\cap D\ne\varnothing\},\notag\\
 Z_D&=\mathcal C(K_D,u|_D,(M[E])_{E\in K_D}).
 \label{eq:promotion-ZD}
\end{align}
There is a direct online equivalence
\begin{equation}
 M[D]\intereq Z_D.
 \label{eq:promotion-restriction-direct}
\end{equation}
If \(D\) meets at least two \(P\)-blocks, the reduced restriction \(M[D]\)
has root partition \(K_D\), actual root preorder \(r|_D\), and child
\(M_B[B\cap D]=M[B\cap D]\) in each cell.  The \(r\)- and \(u\)-comparisons
agree across these actual child colors, so \lemref{lem:colored} changes only
that root.  If \(D\) meets one \(P\)-block, \(Z_D\) is a unary root over
\(M[D]\) and unary suppression proves the same claim.  No child rank is
lifted in either branch, and no recursive hypothesis is used.

Substitute \eqref{eq:promotion-restriction-direct} below the outer \(t\)-root
of \(\operatorname{Ref}(M;R,t)\).  Change that root from \(t\) to \(u\), now
coloring by \(R\), and collapse the two identical \(u\)-selector levels.  The
result is again \(W\), so
\begin{equation}
 \operatorname{Ref}(M;R,t)\intereq W.
 \label{eq:promotion-right-W}
\end{equation}
Together with \eqref{eq:promotion-left-W}, this proves
\eqref{eq:meet-promotion}.  The induction terminates because its sole
recursive call, \eqref{eq:promotion-child-recursion}, is on an actual child
topology.  Constructed refinements such as \(W\) and \(Z_D\) are never used as
induction arguments.
\end{proof}

This factorization also pinpoints why naive height induction fails.  A target
cell may contain types belonging to different actual \(P\)-children, so the
target \(t\)-comparison need not control their actual root comparison.
Meet-promotion uses \(K=P\wedge R\) and a new \(u\) with two safe interfaces:
\(r\) to \(u\) only across \(P\)-colors, and \(t\) to \(u\) only across
\(R\)-colors.  Residues may drift within a color.  Exact parent/inner
occurrence equality is invoked only by the equal-\(u\) tandem law.

\subsection{The main theorem by reduced-topology size}

\begin{theorem}[Topology-recursive flush sufficiency]
\label{thm:topology-recursive-flush-sufficiency}
For every finite alphabet \(A\) and all valid stateless schedulers \(S,T\) on
\(A\),
\[
 S\flusheq T\quad\Longrightarrow\quad S\intereq T.
\]
The proof uses no induction on \(|A|\).
\end{theorem}

\begin{proof}
Prove the statement simultaneously for all finite alphabets by well-founded
recursion on
\begin{equation}
 \nu(S)+\nu(T),
 \label{eq:topology-recursion-measure}
\end{equation}
where both schedulers have first been reduced.  Thus the recursive hypothesis
at a smaller value of \eqref{eq:topology-recursion-measure} is available for
every finite alphabet, not only a smaller subset of the present one.  If
\(A=\varnothing\), every operation is a failed pop, so assume \(A\) nonempty
and write \(F\) for the common flush transformer.

If both reduced schedulers are leaves, let their preorders be
\(\lambda\) and \(\mu\).  The two binary flushes of
\autoref{tab:binary-recovery} give
\[
 \cmp_\lambda(x,y)=\cmp_\mu(x,y)
 \qquad\text{for every \(x,y\in A\)}.
\]
Different types therefore have the same comparison, while occurrences of one
type tie on rank and use the same arrival order.  The two queues couple
occurrence-for-occurrence under every operation word.

Suppose next that \(S\) is a \(\lambda\)-leaf and
\begin{equation}
 T=\mathcal C(Q,s,(T_C)_{C\in Q})
 \label{eq:topology-leaf-internal}
\end{equation}
has an internal root.  Since \(F=\sort_\lambda\), the pair
\((Q,\lambda)\) decomposes \(F\): stable sorting commutes with restriction, and
its root-pattern equation is tautological.  Apply
\lemref{lem:meet-promotion} to both schedulers with this same supplied
decomposition:
\begin{align}
 S&\intereq \mathcal C(Q,\lambda,(S[C])_{C\in Q}),\notag\\
 T&\intereq \mathcal C(Q,\lambda,(T[C])_{C\in Q}).
 \label{eq:topology-leaf-promotions}
\end{align}
For every \(C\in Q\) and word \(w\in C^*\), literal restriction and unary
reduction give
\[
 F_{S[C]}(w)=F(w)=F_{T[C]}(w).
\]
The restriction \(S[C]\) is a leaf, while \(T[C]\) lies strictly below
\(T\)'s root.  Consequently their topology-size sum is strictly smaller than
\eqref{eq:topology-recursion-measure}.  The recursive hypothesis gives
\(S[C]\intereq T[C]\), and child substitution in
\eqref{eq:topology-leaf-promotions} proves \(S\intereq T\).  The case in which
\(T\) is the leaf is symmetric.

Finally suppose both reduced roots are internal:
\begin{equation}
 S=\mathcal C(P,r,(S_B)_{B\in P}),
 \qquad
 T=\mathcal C(Q,s,(T_C)_{C\in Q}).
 \label{eq:topology-both-internal}
\end{equation}
Apply \lemref{lem:finite-meet-decomposition} to their actual decompositions of
\(F\).  It supplies \(R=P\wedge Q\) and \(t\) such that \((R,t)\) decomposes
\(F\).  Meet-promotion gives
\begin{align}
 S&\intereq \mathcal C(R,t,(S[D])_{D\in R}),\notag\\
 T&\intereq \mathcal C(R,t,(T[D])_{D\in R}).
 \label{eq:topology-both-promoted}
\end{align}
Every \(D\in R\) lies in one \(P\)-block and one \(Q\)-block.  Hence both
\(S[D]\) and \(T[D]\) lie strictly below the corresponding original roots,
and for \(w\in D^*\),
\[
 F_{S[D]}(w)=F_S(w)=F_T(w)=F_{T[D]}(w).
\]
Their topology-size sum is strictly smaller than
\eqref{eq:topology-recursion-measure}; the recursive hypothesis yields
\(S[D]\intereq T[D]\).  Substitution into the literally identical outer
\(t\)-selectors in \eqref{eq:topology-both-promoted} completes the proof.

The only recursive arguments in the last two cases are restrictions lying
strictly below at least one original internal root (below both roots in the
internal--internal case).  No constructed refinement is measured, and no
inequality of the form \(|D|<|A|\) is used.  Transporting a finite \(A\) along
a bijection with \(\Fin(|A|)\) recovers \autoref{thm:main} exactly.
\end{proof}

\begin{remark}[Alternative constructions]
\label{rem:alternative-no-induction}
A decomposition--realization route first realizes a finite behavioral
decomposition of one scheduler by structural recursion on its concrete
topology, then builds a pair-dependent common laminar refinement of the two
schedulers.  Its proof deliberately separates measures: raw internal-node count,
including inactive and active-unary nodes, decreases during realization,
whereas the two-scheduler recursion uses the sum of the internal-node counts
of the active, unary-contracted topologies.

A second route is asymmetric: refine \(S\) by the entire topology of \(T\),
flatten the residual \(S\)-subtrees against \(T\)'s leaves, and reconcile
ranks bottom-up.  Its termination measure is lexicographic in the current
\(T\)- and \(S\)-topology sizes, with \(T\) first; because its meet-rank step
uses finiteness, this construction is scoped to finite alphabets and does not
extend verbatim to infinite alphabets.  The extension below is instead a
finite-support reduction and constructs no global infinite intermediate
scheduler.
\end{remark}

\subsection{Arbitrary alphabets by finite support}

The topology-recursive formulation also makes the following finite-support
consequence transparent.

\begin{corollary}[Finite-support extension]
\label{cor:infinite-alphabet-finite-support}
Let \(A\) be an arbitrary type set, finite or infinite, and let \(S,T\) be
valid stateless schedulers with finite topologies.  If their flush outputs
agree on every finite word over \(A\), then their observations agree on every
finite interleaved operation word over \(A\).
\end{corollary}

\begin{proof}
Fix a finite operation word \(o\), and let \(B\subseteq A\) be the finite set
of types pushed in it.  The literal restrictions \(S|_B\) and \(T|_B\) are
valid finite-alphabet schedulers, and their flushes agree on every word over
\(B\) by hypothesis.  Transport \(B\) to \(\Fin(|B|)\) and apply
\autoref{thm:topology-recursive-flush-sufficiency}.  Running \(o\) in either
restriction is exactly its run in the original scheduler because no type
outside \(B\) occurs.  Thus the two observations on \(o\) agree.  Since \(o\)
was arbitrary, the result follows.
\end{proof}

The scope of \corref{cor:infinite-alphabet-finite-support} is deliberately
local to each finite experiment.  The finite meet lemma realizes a finite
constraint graph by natural ranks.  On an infinite alphabet, consistency of
all finite constraint sets does not by itself supply one global
natural-valued rank: an infinite acyclic order can have the wrong order type.
Accordingly the corollary neither constructs nor claims a single
finite-topology canonical form, or a globally chosen meet rank, for the whole
infinite scheduler.  No global canonical form is needed for the observational
conclusion.
